\chapter{Dati e costruzione del campione}
\label{ch:data}

Tutte le fonti impiegate sono pubbliche e scaricabili senza domanda di accesso
né abbonamento. Non è una coincidenza ma un criterio di selezione: mantiene il
progetto entro i tempi di una tesi magistrale e rende i risultati riproducibili
da parte del lettore.

\section{Fonti}
\label{sec:sources}

\begin{table}[H]
\centering
\caption{Fonti dei dati. Tutte pubbliche.}
\label{tab:sources}
\begin{tabularx}{\textwidth}{@{}llX@{}}
\toprule
\textbf{Fonte} & \textbf{Livello} & \textbf{Impiego} \\
\midrule
Eurostat Comext \parencite{eurostat_comext}
  & \acrshort{cn}-8 $\times$ partner $\times$ mese
  & Valore, massa netta e unità supplementari delle importazioni; le variabili
    dipendenti. \\
Regolamento (UE) 2023/956, allegato~I
  & \acrshort{cn}-8
  & Definizione dell'insieme trattato. \\
Serie spot e futures delle \acrshort{eua}
  & Giornaliero
  & Prezzo dei certificati; converte il prelievo in un costo per tonnellata. \\
Parametri settoriali di intensità emissiva
  & Settore $\times$ paese
  & Assegna la dose attesa di trattamento a ciascuna origine. \\
Carbon Pricing Dashboard della Banca Mondiale
  & Paese $\times$ anno
  & Individua le origini il cui prezzo del carbonio è detraibile. \\
\bottomrule
\end{tabularx}
\end{table}

\section{Costruzione di trattamento e controllo}
\label{sec:treatment-control}

L'insieme trattato comprende ogni codice \acrshort{cn}-8 elencato
nell'allegato~I del regolamento. I codici di controllo sono selezionati secondo
una regola fissata prima di osservare qualunque esito, affinché l'insieme di
confronto non risulti scelto in funzione del risultato:

\begin{enumerate}
  \item \textbf{Prossimità.} Il candidato condivide le prime sei cifre del
    \acrshort{hs} con un codice trattato, oppure ricade nella stessa voce a
    quattro cifre laddove non esista un vicino a sei cifre.
  \item \textbf{Esclusione.} Il candidato non compare nell'allegato~I e non è
    previsto in una fase successiva già pubblicata.
  \item \textbf{Non sostituibilità.} Il candidato non è un plausibile sostituto
    dal lato della domanda del suo vicino trattato. È il criterio vincolante e
    quello che richiede maggiore giudizio; la \Cref{sec:substitution} espone
    come venga applicato e come la classificazione risultante sia sottoposta a
    verifica.
  \item \textbf{Supporto.} Entrambi i codici registrano flussi di importazione
    non trascurabili in ogni mese del periodo pre-trattamento, così da
    contribuire all'identificazione e non al rumore.
\end{enumerate}

Il criterio~3 merita particolare enfasi. Un controllo verso cui gli importatori
possano spostarsi è contaminato dal trattamento, e la stima che ne risulta
sovrastima l'effetto contando due volte la sostituzione --- una come calo del
codice trattato e una come aumento del controllo. Quando la sostituibilità è
genuinamente ambigua, la coppia viene esclusa dal campione principale e
reintrodotta in una specifica di robustezza, così che il costo del giudizio
resti visibile.

\section{Periodo campionario}
\label{sec:sample-period}

La finestra di stima va dal gennaio 2021 all'ultimo mese disponibile al momento
della stesura. Si distinguono tre sottoperiodi, e confonderli sarebbe un
errore:

\begin{description}
  \item[Pre-trattamento (2021--set. 2023)] Nessun obbligo \acrshort{cbam} di
    alcun tipo. Fornisce i trend pre-trattamento.
  \item[Transitorio (ott. 2023--dic. 2025)] Sola rendicontazione. Qui è
    possibile un aggiustamento anticipatorio, ed è la ragione per cui questa
    finestra viene stimata con propri coefficienti in tempo d'evento anziché
    accorpata al periodo pre-trattamento.
  \item[Definitivo (gen. 2026--)] Obbligo di certificazione in vigore. È la
    finestra di trattamento.
\end{description}

Trattare la fase transitoria come non trattata attribuirebbe l'eventuale
anticipazione al trend pre-trattamento, distorcendo la stima verso lo zero.
Trattarla come trattata confonderebbe un adempimento documentale con un prezzo.
Viene perciò stimata separatamente, e la differenza fra i due coefficienti è
essa stessa informativa sul fatto che le imprese reagiscano ai costi annunciati
o a quelli effettivamente sostenuti.

\section{Problemi noti dei dati}
\label{sec:data-problems}

\textbf{Soppressione per riservatezza.} Eurostat sopprime le celle che
rivelerebbero un singolo operatore. La soppressione non è casuale --- è
correlata con le origini concentrate, che sono precisamente quelle di
interesse --- sicché le celle soppresse vengono tracciate e riportate anziché
scartate in silenzio.

\textbf{Rumore nei valori unitari.} I \glspl{unitvalue} sono valore diviso
quantità e ne ereditano perciò ogni errore di dichiarazione delle quantità. I
valori estremi sono winsorizzati al primo e al novantanovesimo percentile
entro ciascun codice \acrshort{cn}, e la soglia viene fatta variare nei
controlli di robustezza.

\textbf{Brevità del periodo post-trattamento.} Con due o tre trimestri
osservati, la specifica dinamica dispone di pochi coefficienti in tempo
d'evento successivi al trattamento. Ciò limita quanto si possa affermare sulla
persistenza, e il \Cref{ch:discussion} non afferma più di quanto i dati
sostengano.

\textbf{Lo shock energetico del 2022.} I prezzi dell'energia hanno mosso
violentemente i costi degli input industriali europei nel periodo
pre-trattamento, e lo hanno fatto in modo differenziato fra i settori
energivori --- vale a dire, in modo differenziato fra trattati e controlli. È
la minaccia più seria ai trend paralleli nel campione ed è affrontata
direttamente nella \Cref{sec:threats}.
